%============================================================
% Chapter 10: Minimum Spanning Trees
% STATUS: PLANNED  (generated by tools/outline.py; replace this file when the chapter is written)
%============================================================
\chapter{Minimum Spanning Trees}
\label{ch:mst}

\begin{tcolorbox}[enhanced, colback=white, colframe=secondaryblue, boxrule=1pt, arc=2mm]
\textbf{\large\textcolor{primaryblue}{Connect everything as cheaply as possible --- and understand the cut and cycle properties that make greed correct.}}
\end{tcolorbox}

\chapterfacts{1400}{2200}{5}{10}{30}

\begin{plannedbox}[Planned Contents --- this chapter is not written yet]
\textbf{Part A --- Lessons}
\begin{enumerate}[leftmargin=*, label=\textbf{A\arabic*.}, itemsep=1pt]
    \item Spanning trees; the cut property and the cycle property
    \item Kruskal's algorithm (with DSU)
    \item Prim's algorithm (heap and dense $\bigO{n^2}$ versions); Bor\r{u}vka's algorithm
    \item MST structure: uniqueness, bottleneck paths, second-best MST
    \item Critical edges; MSTs of implicit and geometric graphs; directed MST (minimum arborescence, Chu--Liu/Edmonds)
\end{enumerate}
\textbf{Part B --- Patterns} (with candidate problems, checked when the chapter is written)
\begin{enumerate}[leftmargin=*, label=\textbf{P\arabic*.}, itemsep=1pt]
    \item \textbf{Kruskal and Prim} --- Road Reparation {\small\color{darkgray}(CSES 1675)}, Min Cost to Connect All Points {\small\color{darkgray}(LeetCode 1584)}
    \item \textbf{Virtual vertices} --- Shichikuji and Power Grid {\small\color{darkgray}(CF 1245D)}
    \item \textbf{Which edges can be in an MST?} --- Find Critical and Pseudo-Critical Edges in Minimum Spanning Tree {\small\color{darkgray}(LeetCode 1489)}, Edges in MST {\small\color{darkgray}(CF 160D)}
    \item \textbf{Bottleneck and minimax}
    \item \textbf{Special and directed MSTs} --- 0-1 MST {\small\color{darkgray}(CF 1242B)}, Directed MST {\small\color{darkgray}(Library Checker)}
\end{enumerate}
\end{plannedbox}
